\section{Deep Q-Learning}
    \subsection{Main Idea}
        The \textbf{Deep Q-Learning} \cite{Mnih2013} is one of the earliest and successful deep reinforcement learning algorithms.
        Suprisingly, DQN algorithm are not guaranteed to converge, since the deadly triad perfectly fits into the algorithms' condition:
        it bootstraps, it is off-policy, and it uses function approximation. But there are practical methods for calming down this unstability.
        But still note that the algorithm is not guaranteed to converge.

    \subsection{Objective Function}
        The objective function which DQN minimizes is given as following:
        \begin{equation}
            J = \E\Big[\Big( R + \gamma \max_{a \in \mathcal{A}(S')}{\hat{q}(S',a,w)} - \hat{q}(S, A, w) \Big)^2\Big],
        \end{equation}
        where $\hat{q}$ is the approximation function of action values
        and $(S, A, R, S')$ are random variables denoting state, action, reward and next state, resepectively.
        Note that the squared term is a Bellman optimality error.

        Now, we can caculate the gradient and minimize the loss with gradient descents... but take a look. Unfortunately, the
        max operation is not differentiable. To resolve this, we freeze the parameter of a network inside the max operation for a short time. Specifically,
        we introduce two networks: \emph{target network and main network}. Using the main network $\hat{q}(s, a, w)$ and $\hat{q}(s, a, w_T)$,
        we can re-write the objective function as following:
        \begin{equation}
            J = \E\Big[\Big( R + \gamma \max_{a \in \mathcal{A}(S')}{\hat{q}(S',a,w_T)} - \hat{q}(S, A, w) \Big)^2\Big], \label{eq:DQNObjective}
        \end{equation}
        and the gradient can be computed about $w$. The parameters of target network and main network are synchronized regularly.

        There is one more technique required for DQN to be trained. The \emph{experience replay} collects the $(s, a, r, s')$ pair, and then
        samples the experiences with uniform random distrubution. Why is this technique required?
        
        \begin{itemize}
            \item \cite{Zhao2025} says: Compare equation \ref{eq:DQNObjective} and \ref{eq:stationaryDist}. Since we don't know the probability distribution of $(S, A, R, S')$,
                we just simply set it as uniform distrubution. However, the state-action pairs are are strongly correlated, and not uniformly distributed in single trajectory.
                To break the correlation, we uniformly draw samples from the replay buffer.
        \end{itemize}
        
        Also, since we can reuse the experiences, the sample efficiency naturally increases.
    \subsection{Pseudocode}
        \Algorithm{DQN}
        {Deep Q-Learning}
        {
            \inputitem{$b$}{behavior policy} \\
            \inputitem{$C$}{sync rate}
        }
        {\inputitem{$\hat{q}(s, a, w)$}{approximated action values}}
        {
            store the samples generated by $b$ in a replay buffer $\mathcal{B} = \{(s, a, r, s')\}$\;
            \ForEach{iteration}
            {
                uniformly draw a mini batch $B$ of samples from $\mathcal{B}$\;
                $L \leftarrow 0$\;
                \ForEach{sample $(s, a, r, s') \in B$}
                {
                    $y_T \leftarrow r + \gamma \max_{a \in \mathcal{A}(s')}{\hat{q}(s', a, w_T)}$\;
                    $L \leftarrow L + \Big( y_T - \hat{q}(s, a, w) \Big)^2$\;
                }
                minimize $L$ by updating parameter $w$\;
                set $w_t = w$ every $C$ iterations
            }
            \Return{$\hat{q}$}
            
        }
    \subsection{Implementation Details}
        Since the original DQN algorithm is highly unstable, there are several techniques to stabilize the algorithm.
        \begin{itemize}
            \item Priortized experience replay (PER)
            \item Double DQN
            \item Dueling DQN
            \item Distributional DQN
        \end{itemize}
    \subsection{Pros and Cons}
        Pros:
        \begin{itemize}
            \item 
        \end{itemize}
        Cons:
        \begin{itemize}
            \item Catastrophic forgetting: This is a common problems in neural networks.
            \item No convergence guaranteed: Note that DQN satisfies the deadly triad condition.
                And also, the algorithm does not even use the true gradient, since the target network is freezed for a moment. This is called a \textbf{semi-gradient} method,
                and no convergence is guaranteed for this method. For more information see \cite{Sutton2020}.
        \end{itemize}